
\chapter{Introduction}

\section{Motivation}

Data-to-text (D2T) generation concerns the automatic conversion of structured
information into natural-language text. It is useful whenever information stored
as records, tables, databases, time series, or knowledge graphs must be converted
in a linguistic form, which is more suitable for a human reader, a conversational
interface, or another language-based system. D2T systems must therefore solve
two related problems: they must select and organize the information that should
be communicated, and then express that information in text that is fluent,
coherent, and faithful to its source
\cite{gatt2018survey,osuji2024systematic}.

Recent neural and large language model based systems have substantially improved
the fluency and flexibility of generated text. However, fluent text does not
necessarily preserve the content of the input. A system may omit a fact,
introduce an unsupported statement, select different facts on different
executions in a way that changes their apparent
meaning. Additionally, a direct neural generator, these intermediate choices are not usually
represented as explicit objects that a developer can inspect, modify, or compare
with the source. This makes it difficult to identify whether an error resulted
from content selection, content ordering, aggregation, lexicalization, or surface
realization.
\cite{wiseman-etal-2017-challenges,dusek-kasner-2020-evaluating,
dhingra-etal-2019-handling}.

Before neural generation became dominant, D2T systems were commonly built using
manually authored rules and templates. A rule-based system specifies conditions
and operations for selecting, ordering, and combining facts, while a template
provides a sentence pattern into which values can be inserted. Such systems are
easy to inspect and can provide strong control over which source facts are
realized. Their limitation is not that rules are inherently expensive, but that
the engineering effort grows as the task becomes more varied. New predicates,
combinations of facts, aggregation patterns, domains, and linguistic constructions
may each require additional design, testing, and maintenance
\cite{reiter1997building,gatt2018survey}.

This challenge is especially visible for complex structured data. A time series,
for example, may require decisions about temporal ordering, aggregation,
extremes, trends, and the level of detail appropriate for the output. These
decisions are not caused by the presence of time alone; they arise because the
system must handle many possible structures and realizations. Neural systems
reduce some of this manual engineering and offer greater linguistic variety, but
their learned generation behavior is harder to inspect.

One way to make this trade-off more manageable is to separate meaning
representation from surface realization. A meaning representation is an
intermediate, structured description of the content to be communicated. It is
independent of the exact wording of the final text and can be expressed using,
for example, logical forms, semantic frames, text plans, or graphs. An explicit
meaning representation makes the selected content available for inspection before
it is verbalized. The resulting decomposition can be summarized as
\begin{equation}
    d \longrightarrow T \longrightarrow y,
\end{equation}
where $d$ denotes a structured data instance, $T$ denotes its interpretable meaning
representation, and $y$ denotes the generated text. In this context,
interpretability means that the intermediate meaning representation and the procedure
that produces the text are explicit artifacts that can be inspected.
Meaning representation make the facts being communicated explicit and provide an intermediate
object that can be inspected, normalized, and evaluated independently of the
final wording. Nevertheless, constructing consistent meaning representation from complex data
and producing fluent text from it remains non-trivial tasks.

Existing work has shown that LLMs can be used to construct an interpretable
rule-based meaning representation to text generator.
Lango and Dušek~\cite{lango-dusek-2025-llm} use cooperating LLM agents to design
and refine such a generator (a Python program), from meaning representation, which can be inspected, executed without an LLM at inference time, and corrected by a developer when necessary.

The scope of that approach begins after a meaning representation has already
been provided: it addresses just the surface realization.
A further limitation of a fixed sequence of agent interactions is that it does
not by itself define a systematic criterion for comparing many alternative
generators. Agents can propose and critique changes, but their next
action is primarily guided by the current conversation.

For these reasons, the objective of this thesis is to considers the complete D2T pipeline, from heterogeneous
raw structured instances through an explicit meaning representation to the final
text.
Separating these two stages also
allows errors introduced during meaning-representation construction to be
distinguished from errors introduced during surface realization.
In the first stage an
LLM-based pipeline is used for the meaning-representation construction. A domain-specialized
model is also examined as a way of producing both the meaning representation and a text in a single inference pass.
In the second stage, the surface-realization method used in this thesis places an
LLM inside an evolutionary optimization loop inspired by AlphaEvolve
\cite{novikov2025alphaevolve}, a framework for evolutionary program improvement
using LLM-generated mutations and automatic evaluation.
The advantage of such an approach over the agents is that it evaluates every proposed program using the same
evaluator, retains high-scoring and structurally diverse candidates, and uses
the accumulated results to guide subsequent modifications. This makes it
possible to reject plausible but ineffective changes and to revisit alternative
realization strategies rather than relying on a single refinement trajectory.
Reference-based and semantic metrics are used to evaluate generation quality and faithfulness.
Interpretability is supported by exposing the meaning representation and the human-readable program code.

The work makes the following contributions:
\begin{itemize}
    \item three LLM-based pipeline designs are proposed for constructing meaning representation
    from heterogeneous structured data;
    \item a domain-specialized training procedure that distills the
    conversion process into a large language model;
    \item an evolutionary surface-realization method in which LLM-generated
    program mutations are executed and selected using automatic evaluation;
    \item a systematic empirical evaluation of the proposed methods across
    multiple domains using reference-based, semantic, and LLM-based measures.
\end{itemize}

\section{Thesis structure}

The structure of the thesis is as follows. Chapter~2 introduces the D2T task,
reviews the principal generation approaches, discusses interpretability, and
describes the evaluation methods used in this area. Chapter~3 presents the
meaning-representation construction stage of the proposed pipeline. It describes
how complex structured instances are converted into meaning representation using
LLM-based pipelines and a domain-specialized fine-tuned model.

Chapter~4 presents the surface-realization stage, namely the conversion of
semantic triples into natural-language text. It introduces the AlphaEvolve and
OpenEvolve frameworks and explains how an executable program for verbalizing
triples is selected, mutated, evaluated, and retained during evolution.
Chapter~5 describes the experimental datasets and configurations, reports the
results for meaning-representation construction and surface realization, and
evaluates the complete pipeline. Chapter~6 discusses the results, limitations,
and broader implications of the proposed approach. Finally, Chapter~7
summarizes the thesis and presents its conclusions. The appendices contain the
prompts used by the conversion methods and the initial surface-realization
program.
